% frontmatter/slowa/como-usar.tex -- la página para el adulto de
% «Pierwsze słowa»: la de frontmatter/palabras/como-usar.tex, en polaco,
% con lo que cambia en polaco (las 32 letras, las que se leen distinto
% que en español).
\section*{Jak korzystać z tego zeszytu}

To pierwszy zeszyt z serii \emph{Uczę się czytać} po polsku. Jest dla
dziecka, które zna już litery (albo prawie wszystkie) i zaczyna je
łączyć: czyta sylaby i chce przejść od sylab do całych słów. Na końcu
zeszytu będzie czytać krótkie zdania.

Na każdy dzień roku od poniedziałku do piątku jest jedna strona, także
w wakacje. Wystarczy pięć albo dziesięć minut dziennie. Każda strona
wygląda tak samo:

\begin{itemize}[leftmargin=6mm]
\item \textbf{Data}, \textbf{Czytanie} i \textbf{Uwagi}: dla
  dorosłego. Zaznaczenie codziennie, czy dziecko czytało samodzielnie,
  z pomocą czy z trudem, nic nie kosztuje, a po roku jest zapisem
  wszystkich postępów.
\item \textbf{Słowo dnia}, w zielonej ramce: u góry podzielone na
  sylaby (\emph{ko · ło}), pod spodem całe (\emph{koło}). Najpierw
  czyta się sylaba po sylabie, powoli, a potem całe słowo, płynnie. Na
  początku jest jedno słowo dziennie; zimą dwa; wiosną trzy; latem
  krótkie zdanie.
\item \textbf{Zadanie}, które prawie zawsze kończy się rysowaniem.
  Polecenia (małą szarą czcionką) \textbf{czyta dorosły} na głos:
  dziecko czyta tylko to, co jest napisane dużymi literami.
\end{itemize}

\subsection*{Zadania}

{\small
\begin{itemize}[leftmargin=6mm, itemsep=1pt, topsep=2pt]
\item \textbf{\lblDibuja}: narysować to, co się właśnie przeczytało.
  To najprostszy sposób, żeby sprawdzić, że słowo zostało zrozumiane, a
  nie tylko odczytane.
\item \textbf{\lblCompleta}: kilka linii bez kształtu, z których trzeba
  zrobić rysunek (buzię taty, łódkę, budę Toby'ego...).
\item \textbf{\lblTraza}: poprowadzić palcem albo ołówkiem po literze,
  wielkiej i małej -- pierwszej literze słowa dnia -- a potem napisać ją
  samodzielnie na linii. W ciągu roku pojawiają się wszystkie 32 litery
  polskiego alfabetu; ą, ę, ń i y, od których nie zaczyna się żadne
  polskie słowo, ze słowem, w którym są w środku (\emph{wąż},
  \emph{koń}).
\item \textbf{\lblEscribe}: obwieść słowo dnia w środku jego pustych
  liter, a potem napisać je bez pomocy.
\item \textbf{\lblEncuentra}: zakreślić słowo dnia za każdym razem, kiedy
  pojawia się między innymi, bardzo do niego podobnymi (\emph{łapa,
  łata, mapa...}). Trzeba patrzeć na wszystkie litery, nie tylko na
  pierwszą.
\item \textbf{\lblPalmadas}: czytać słowo, klaszcząc przy każdej
  sylabie, i pokolorować jedno kółko za każde klaśnięcie.
\item \textbf{\lblAdivina}: dorosły czyta zagadkę, a dziecko rysuje
  odpowiedź (od zimy pisze też pod spodem, co to jest). Odpowiedzią jest
  zwykle słowo przeczytane w tym tygodniu.
\item \textbf{\lblUne}: połączyć linią każde słowo z tym samym słowem
  WIELKIMI LITERAMI -- latem każde zwierzę z jego odgłosem, każde słowo z
  przeciwnym albo z tym, które się z nim rymuje.
\item \textbf{\lblSiNo} (tylko latem): czytać bardzo krótkie zdania i
  decydować, czy mają sens (\emph{Toby szczeka} -- tak; \emph{Toby
  lata} -- nie).
\item \textbf{\lblRelee} i \textbf{\lblRepasa} (w piątki): wszystkie
  słowa z tygodnia razem, żeby przeczytać je jeszcze raz i zaznaczyć te,
  które wychodzą już same. Co dziesięć stron jest mały napis, żeby to
  uczcić.
\end{itemize}}

\subsection*{Jak pomagać}

{\small
\begin{itemize}[leftmargin=6mm, itemsep=1pt, topsep=2pt]
\item Jeśli dziecko utknie na jakimś słowie, warto dać mu kilka sekund,
  a jeśli nie wychodzi, przeczytać je, pokazując palcem sylaby -- i
  poprosić, żeby je powtórzyło. Lepsze to niż zamieniać stronę w walkę.
\item Nie trzeba poprawiać każdego błędu. Jeśli przeczytało \emph{koza}
  zamiast \emph{kosa}, wystarczy zapytać: „na pewno? popatrz jeszcze
  raz”.
\item Jeśli przez cały tydzień czyta słowo dnia od razu, nie trzeba
  czekać: można czytać więcej niż jedną stronę dziennie. Jeśli idzie
  bardzo ciężko, można powtórzyć stronę z poprzedniego dnia.
\item Zadania nie trzeba kończyć tego samego dnia. Najważniejsze jest
  czytanie; rysunek może poczekać do popołudnia.
\item Jeśli dziecko uczy się czytać także po hiszpańsku, jak Lucía,
  warto na początku przećwiczyć litery, które po polsku czyta się
  inaczej: \emph{c} (\emph{taca}), \emph{j} (\emph{jaja}), \emph{h}
  (\emph{hak}), \emph{w} (\emph{woda}), \emph{z} (\emph{koza}),
  \emph{y} (\emph{ryby}), \emph{ł} (\emph{łapa}), a zimą \emph{ó}, które
  czyta się jak \emph{u}.
\end{itemize}}
\newpage
